\documentclass[11pt]{article}
% ===========================================================================
%  Shared preamble for the AP Calculus AB exam-preparation worksheets.
%  Input right after \documentclass. Files beginning with "_" are not compiled
%  on their own. Navy and gold, to match the course's lesson pages.
% ===========================================================================
\usepackage[letterpaper,top=2.3cm,bottom=2.2cm,left=1.7cm,right=1.7cm,headheight=28pt,footskip=30pt]{geometry}
\usepackage{amsmath,amssymb}
\usepackage[T1]{fontenc}
\usepackage[utf8]{inputenc}
\usepackage{lmodern}
\usepackage{xcolor}
\usepackage[most]{tcolorbox}
\usepackage{enumitem}
\usepackage{fancyhdr}
\usepackage{array}
\usepackage{tikz}
\usepackage{pgfplots}
\pgfplotsset{compat=1.18}

\definecolor{navy}{HTML}{161E54}
\definecolor{ink}{HTML}{27357E}
\definecolor{gold}{HTML}{B98A00}
\definecolor{goldbg}{HTML}{FFFDF3}
\definecolor{inkbg}{HTML}{F6F7FD}
\definecolor{rule}{HTML}{C9CFF0}
\definecolor{okgreen}{HTML}{1B7A44}
\definecolor{okbg}{HTML}{F1F8F3}

\pagestyle{fancy}
\fancyhf{}
\fancyhead[L]{\footnotesize Name: \underline{\hspace{4cm}}}
\fancyhead[C]{\textbf{\color{navy}Integration Academy}\\[-2pt]{\scriptsize Calculus AB \textperiodcentered{} AP\textsuperscript{\textregistered} Exam Preparation}}
\fancyfoot[L]{\tiny\color{black!55}\parbox[t]{0.78\textwidth}{AP\textsuperscript{\textregistered} is a registered trademark of the College Board, which was not involved in the production of, and does not endorse, this material. \textcopyright{} Integration Academy.}}
\fancyfoot[R]{\scriptsize\thepage}
\renewcommand{\headrulewidth}{1.2pt}
\renewcommand{\headrule}{\hbox to\headwidth{\color{gold}\leaders\hrule height \headrulewidth\hfill}}
\renewcommand{\footrulewidth}{0pt}
\setlength{\parindent}{0pt}
\setlength{\parskip}{4pt}

% Title block: \wstitle{code}{title}{unit line}
\newcommand{\wstitle}[3]{\begin{center}{\small\bfseries\color{gold}WORKSHEET #1}\\[3pt]{\LARGE\bfseries\color{navy}#2}\\[3pt]{\small\color{black!60}#3}\end{center}\vspace{2pt}}
% Section heading
\newcommand{\sech}[1]{\par\vspace{10pt}{\large\bfseries\color{navy}#1}\par\nobreak\vspace{2pt}{\color{rule}\hrule height0.8pt}\vspace{7pt}}
% A level tag printed beside a group of questions
\newcommand{\level}[1]{\par\vspace{6pt}{\footnotesize\bfseries\color{gold}\MakeUppercase{#1}}\par\nobreak\vspace{2pt}}

\newtcolorbox{ideabox}{colframe=ink,colback=inkbg,coltitle=white,colbacktitle=ink,fonttitle=\bfseries\small,title={Key ideas},arc=2pt,boxrule=0.8pt,left=8pt,right=8pt}
\newtcolorbox{example}[1]{colframe=rule,colback=white,coltitle=navy,colbacktitle=inkbg,fonttitle=\bfseries\small,title={#1},arc=2pt,boxrule=0.8pt,borderline west={3pt}{0pt}{ink},left=10pt,right=8pt}
\newtcolorbox{frbox}[1]{colframe=gold,colback=goldbg,coltitle=white,colbacktitle=gold,fonttitle=\bfseries\small,title={#1},arc=2pt,boxrule=0.8pt,breakable,left=8pt,right=8pt}
\newtcolorbox{solbox}{colframe=okgreen,colback=okbg,coltitle=white,colbacktitle=okgreen,fonttitle=\bfseries\small,title={Worked solutions},arc=2pt,boxrule=0.8pt,breakable,left=8pt,right=8pt}
\newtcolorbox{scorebox}[1]{colframe=navy,colback=white,coltitle=white,colbacktitle=navy,fonttitle=\bfseries\small,title={#1},arc=2pt,boxrule=0.8pt,breakable,left=8pt,right=8pt}

\newcommand{\soln}{\par\smallskip\textbf{\color{ink}Solution.}\ }
\newcommand{\result}[1]{\par\smallskip\textbf{#1}}
% blank working space of a given height
\newcommand{\work}[1]{\par\vspace{#1}}
% a point value in the margin of a scoring guide
\newcommand{\pts}[1]{\hfill{\footnotesize\bfseries\color{gold}[#1]}}

% Coordinate grid for graph questions: \begin{apgraph}{xmin}{xmax}{ymin}{ymax} ... \end{apgraph}
\newenvironment{apgraph}[4]{\begin{center}\begin{tikzpicture}\begin{axis}[
  axis lines=middle, xmin=#1, xmax=#2, ymin=#3, ymax=#4,
  grid=both, grid style={black!12}, xtick distance=1, ytick distance=1,
  ticklabel style={font=\scriptsize}, xlabel={$x$}, ylabel={$y$},
  width=10cm, height=7.2cm, samples=80, clip=false,
  every axis plot/.append style={line width=1.1pt, color=ink}]}{\end{axis}\end{tikzpicture}\end{center}}
\newcommand{\closeddot}[2]{\addplot[only marks, mark=*, mark size=2.2pt, color=ink] coordinates {(#1,#2)};}
\newcommand{\opendot}[2]{\addplot[only marks, mark=*, mark size=2.2pt, color=ink, mark options={fill=white}] coordinates {(#1,#2)};}

% Multiple-choice options: four across, or stacked when the options are long.
\newcommand{\opts}[4]{\par\vspace{3pt}\begin{tabular}{@{}p{0.235\linewidth}p{0.235\linewidth}p{0.235\linewidth}p{0.235\linewidth}@{}}(A)\ #1 & (B)\ #2 & (C)\ #3 & (D)\ #4\end{tabular}\par\vspace{9pt}}
\newcommand{\optsv}[4]{\par\vspace{3pt}\begin{tabular}{@{}l@{\ \ }p{0.9\linewidth}@{}}(A) & #1\\(B) & #2\\(C) & #3\\(D) & #4\end{tabular}\par\vspace{9pt}}

\fancyhead[R]{\footnotesize\color{black!60}Worksheet 6.3}
\begin{document}

\wstitle{6.3}{From Sums to the Definite Integral}{Unit 6: Integration and Accumulation of Change}

\begin{ideabox}
\begin{itemize}[leftmargin=*,itemsep=2pt,topsep=2pt]
  \item $\displaystyle\int_{a}^{b}f(x)\,dx=\displaystyle\lim_{n\to\infty}\sum_{i=1}^{n}f(x_i)\,\Delta x$, with $\Delta x=\dfrac{b-a}{n}$ and $x_i=a+i\,\Delta x$.
  \item To read a sum as an integral: the factor $\dfrac{b-a}{n}$ gives the width of the interval, and the expression $a+i\,\Delta x$ gives where it starts.
  \item A definite integral is a \textbf{signed area}: regions below the axis count as negative.
  \item Integrals of linear functions and semicircles can be evaluated from geometry, with no antiderivative.
\end{itemize}
\end{ideabox}

\sech{Worked examples}
\begin{example}{Example 1 \textperiodcentered{} A sum as an integral}
Write $\displaystyle\lim_{n\to\infty}\sum_{i=1}^{n}\left(1+\frac{3i}{n}\right)^2\cdot\frac{3}{n}$ as a definite integral.\soln $\Delta x=\dfrac3n$, so $b-a=3$; $x_i=1+\dfrac{3i}{n}$, so $a=1$ and $b=4$.\result{$\displaystyle\int_{1}^{4}x^2\,dx$.}
\end{example}

\begin{example}{Example 2 \textperiodcentered{} A semicircle}
Evaluate $\displaystyle\int_{-2}^{2}\sqrt{4-x^2}\,dx$.\soln The graph is the upper half of a circle of radius 2.\result{$\dfrac12\pi(2)^2=2\pi$.}
\end{example}

\begin{example}{Example 3 \textperiodcentered{} An absolute value}
Evaluate $\displaystyle\int_{0}^{5}|x-3|\,dx$.\soln Two triangles above the axis: area $\dfrac12(3)(3)$ on $[0,3]$ and $\dfrac12(2)(2)$ on $[3,5]$.\result{$\dfrac92+2=\dfrac{13}{2}$.}
\end{example}

\newpage
\sech{Practice}

\level{Core}
\begin{enumerate}[leftmargin=*,itemsep=2pt]
  \item Evaluate $\displaystyle\int_{0}^{6}3\,dx$ using geometry.\work{2.3cm}
  \item Evaluate $\displaystyle\int_{0}^{4}x\,dx$ using geometry.\work{2.3cm}
  \item Evaluate $\displaystyle\int_{-1}^{1}\sqrt{1-x^2}\,dx$ using geometry.\work{2.3cm}
  \item Write a definite integral for the area under $y=x^3$ from $x=1$ to $x=2$.\work{2.3cm}
\end{enumerate}

\level{Exam level}
\begin{enumerate}[leftmargin=*,itemsep=2pt,start=5]
  \item Write $\displaystyle\lim_{n\to\infty}\sum_{i=1}^{n}\cos\!\left(\frac{3i}{n}\right)\cdot\frac{3}{n}$ as a definite integral.\work{2.8cm}
  \item Write $\displaystyle\int_{1}^{3}\dfrac1x\,dx$ as the limit of a right Riemann sum.\work{2.8cm}
  \item Evaluate $\displaystyle\int_{-2}^{4}(2-x)\,dx$ using geometry.\work{2.8cm}
  \item Evaluate $\displaystyle\int_{0}^{3}\left(1+\sqrt{9-x^2}\right)\,dx$ using geometry.\work{2.8cm}
\end{enumerate}

\begin{frbox}{Question 9 \textperiodcentered{} Free-response practice}
Let $g(x)=2x-4$.
\begin{enumerate}[label=(\alph*),leftmargin=*,itemsep=2pt]
  \item Find the areas of the two triangles between the graph of $g$ and the $x$-axis for $0\le x\le5$.\work{2.6cm}
  \item Use your answer to evaluate $\displaystyle\int_{0}^{5}g(x)\,dx$.\work{2.6cm}
  \item Evaluate $\displaystyle\int_{0}^{5}|g(x)|\,dx$, and explain why it differs from part (b).\work{2.6cm}
\end{enumerate}
\end{frbox}

\newpage
\begin{solbox}
\begin{enumerate}[leftmargin=*,itemsep=4pt]
  \item A rectangle: 18.
  \item A triangle: $\dfrac12(4)(4)=8$.
  \item A semicircle of radius 1: $\dfrac\pi2$.
  \item $\displaystyle\int_{1}^{2}x^3\,dx$.
  \item $\Delta x=\dfrac3n$ and $x_i=\dfrac{3i}{n}$: $\displaystyle\int_{0}^{3}\cos x\,dx$.
  \item $\Delta x=\dfrac2n$, $x_i=1+\dfrac{2i}{n}$: $\displaystyle\lim_{n\to\infty}\sum_{i=1}^{n}\dfrac{1}{1+\frac{2i}{n}}\cdot\dfrac2n$.
  \item Above the axis on $[-2,2]$: $\dfrac12(4)(4)=8$. Below on $[2,4]$: $\dfrac12(2)(2)=2$. Value $8-2=6$.
  \item A rectangle of area 3 plus a quarter circle of radius 3: $3+\dfrac{9\pi}{4}$.
  \item \begin{enumerate}[label=(\alph*),leftmargin=*,itemsep=3pt]
    \item Below the axis on $[0,2]$: $\dfrac12(2)(4)=4$. Above on $[2,5]$: $\dfrac12(3)(6)=9$.
    \item $-4+9=5$.
    \item $4+9=13$. The absolute value counts the region below the axis as positive.
  \end{enumerate}
\end{enumerate}
\end{solbox}

\end{document}
